\ifdefined\gkver\else\def\gkver{student}\fi
\documentclass[\gkver]{gaokaozhenti}
\begin{document}

\chapter{2017年江苏}

\begin{enumerate}
\item
%% number: 1
%% paperName: 2017年江苏高考第 1 题 3 分
%% typeId: 1
%% type: 单选题
%% chapter: 磁场
%% point: 磁通量
%% method: 无
%% score: 3
%% degree: 750
%% duplicateId: 0
%% body:
如图所示，两个单匝线圈 a、b 的半径分别为 $r$ 和 $2r$。圆形匀强磁场 $B$ 的边缘恰好与 a 线圈重合，则穿过 a、b 两线圈的磁通量之比为\xzanswer{A}

\onepicture[w=4.2cm]{\includesvg{figs/01}}

\fourchoices[answer=A]
{$1:1$}
{$1:2$}
{$1:4$}
{$4:1$}

%% answer: A
%% memo:
\memoanswer{
由磁通量的定义式 $\Phi = BS$，其中 $S$ 为垂直磁感线方向的有效面积，结合题图可知，穿过 $a$、$b$ 两个线圈的磁通量均为 $\Phi = B\cdot\pi r^{2}$，因此穿过 $a$、$b$ 两个线圈的磁通量比值为 1:1，故 A 正确。
}

\item
%% number: 2
%% paperName: 2017年江苏高考第 2 题 3 分
%% typeId: 1
%% type: 单选题
%% chapter: 曲线运动
%% point: 平抛运动
%% method: 无
%% score: 3
%% degree: 650
%% duplicateId: 0
%% body:
如图所示，A、B 两小球从相同高度同时水平抛出，经过时间 $t$ 在空中相遇，若两球的抛出速度都变为原来的 2 倍，则两球从抛出到相遇经过的时间为\xzanswer{C}

\onepicture[w=5.2cm]{\includesvg{figs/02}}

\fourchoices[answer=C]
{$t$}
{$\frac{\sqrt 2}{2}t$}
{$\frac{t}{2}$}
{$\frac{t}{4}$}

%% answer: C
%% memo:
\memoanswer{
设两球开始时水平距离为 $L$，两球第一次抛出的速度大小分别为 $v_{1}$、$v_{2}$，由于小球抛出后水平方向做匀速直线运动，则从抛出到两球相遇的时间为 $t = \frac{L}{v_{1} + v_{2}}$，若两个球的抛出速度都变为原来的 2 倍，则从抛出到相遇经过的时间为 $t' = \frac{L}{2(v_{1} + v_{2})} = \frac{t}{2}$，故 C 正确。
}

\item
%% number: 3
%% paperName: 2017年江苏高考第 3 题 3 分
%% typeId: 1
%% type: 单选题
%% chapter: 功和能
%% point: 功能中的图象
%% method: 无
%% score: 3
%% degree: 550
%% duplicateId: 0
%% body:
一小物块沿斜面向上滑动，然后滑回到原处。物块初动能为 $E_{k0}$，与斜面间的动摩擦因数不变，则该过程中，物块的动能 $E_{k}$ 与位移 $x$ 关系的图线是\xzanswer{C}

\fourchoices[answer=C, ispicture=true, h=2.4cm]
{figs/03a.png}
{figs/03b.png}
{figs/03c.png}
{figs/03d.png}

%% answer: C
%% memo:
\memoanswer{
设物块与斜面间的动摩擦因数为 $\mu$，物块的质量为 $m$，重力加速度为 $g$，斜面的倾角为 $\theta$，则物块在上滑过程中，由动能定理得到动能 $E_{k}$ 与位移 $x$ 的关系为 $-(mg\sin\theta + \mu mg\cos\theta)x = E_{k} - E_{k0}$，即 $E_{k} = E_{k0} - (mg\sin\theta + \mu mg\cos\theta)x$，物块沿斜面下滑时，同样由动能定理得 $(mg\sin\theta - \mu mg\cos\theta)(x_{0} - x) = E_{k}$，此过程动能 $E_{k}$ 与位移 $x$ 的关系为 $E_{k} = (mg\sin\theta - \mu mg\cos\theta)(x_{0} - x)$，由此可以判断，故 C 正确。
}

\item
%% number: 4
%% paperName: 2017年江苏高考第 4 题 3 分
%% typeId: 1
%% type: 单选题
%% chapter: 电场
%% point: 带电粒子的直线运动
%% method: 无
%% score: 3
%% degree: 550
%% duplicateId: 0
%% body:
如图所示，三块平行放置的带电金属薄板 A、B、C 中央各有一小孔，小孔分别位于 O、M、P 点。由 O 点静止释放的电子恰好能运动到 P 点。现将 C 板向右平移到 $P'$ 点，则由 O 点静止释放的电子\xzanswer{A}

\onepicture[w=5.4cm]{\includesvg{figs/04}}

\fourchoices[answer=A]
{运动到 P 点返回}
{运动到 P 和 $P'$ 点之间返回}
{运动到 $P'$ 点返回}
{穿过 $P'$ 点}

%% answer: A
%% memo:
\memoanswer{
粒子经 A、B 板间的电场加速，在 B、C 板间的电场中减速，设 A、B 板间的电压为 $U$，B、C 板间的电场强度大小为 $E$，M、P 间的距离为 $d$，由动能定理得 $eU - eEd = 0$；若将 C 板向右平移到 $P'$ 点，由于 B、C 两板的带电量一定，由 $E = \frac{U}{d} = \frac{Q}{Cd} = \frac{4\pi kQ}{\varepsilon S}$ 可知，C 板向右移动过程中 B、C 两板间的电场强度不变，C 板移到 $P'$ 点后，电子仍由 O 点静止释放，设在 B、C 板间向右运动的最大距离为 $d'$，同样由动能定理有 $eU - eEd' = 0$，求得 $d' = d$，由此可以判断，电子经 A、B 加速后在 B、C 间仍减速到 P 点时速度为零，然后返回，故 A 正确 BCD 错误。
}

\item
%% number: 5
%% paperName: 2017年江苏高考第 5 题 3 分
%% typeId: 1
%% type: 单选题
%% chapter: 曲线运动
%% point: 竖直平面圆周运动
%% method: 临界
%% score: 3
%% degree: 350
%% duplicateId: 0
%% body:
如图所示，一小物块被夹子夹紧，夹子通过轻绳悬挂在小环上，小环套在水平光滑细杆上，物块质量为 $M$，到小环的距离为 $L$，其两侧面与夹子间的最大静摩擦力均为 $F$。小环和物块以速度 $v$ 向右匀速运动，小环碰到杆上的钉子 P 后立刻停止，物块向上摆动。整个过程中，物块在夹子中没有滑动。小环和夹子的质量均不计，重力加速度为 $g$。下列说法正确的是\xzanswer{D}

\onepicture[w=4.2cm]{\includesvg{figs/05}}

\fourchoices[answer=D]
{物块向右匀速运动时，绳中的张力等于 $2F$}
{小环碰到钉子 P 时，绳中的张力大于 $2F$}
{物块上升的最大高度为 $\frac{2v^{2}}{g}$}
{速度 $v$ 不能超过 $\sqrt{\frac{(2F - Mg)L}{M}}$}

%% answer: D
%% memo:
\memoanswer{
由于不计夹子的质量，因此绳的张力大小始终等于夹子和物块间的摩擦力，物块向右匀速运动时，由物体的平衡知，绳中的张力等于物块的重力 $Mg$，由于物块向上摆动时做圆周运动，因此刚开始摆动时，夹子对物块的摩擦力仍为静摩擦力且大于 $Mg$，此时的摩擦力小于等于 $2F$，由此可以判断，匀速运动时绳中张力小于 $2F$，故 A 错误；当小环碰到钉子 P 时，夹子与物块间的摩擦力仍为静摩擦力，由静摩擦力与物块重力的合力提供向心力，此时夹子没有移动，则此时绳中的张力仍小于等于 $2F$，故 B 错误；如果物块上升的最大高度不超过细杆，由机械能守恒可知，$Mgh = \frac{1}{2}Mv^{2}$，即上升的最大高度 $h = \frac{v^{2}}{2g}$，当物块上摆的高度超过细杆时，若物块不能做完整的圆周运动，则过细杆后物块做斜抛运动，在最高点速度水平且不为 0，则上升高度小于 $\frac{v^{2}}{2g}$，故 C 错误；当物块向上摆动的一瞬间，如果物块与夹子间的静摩擦力刚好为 $2F$，这时的 $v$ 是允许的最大速度，则 $2F - Mg = M\frac{v^{2}}{L}$，求得 $v = \sqrt{\frac{(2F - Mg)L}{M}}$，故 D 正确。
}

\item
%% number: 6
%% paperName: 2017年江苏高考第 6 题 4 分
%% typeId: 2
%% type: 多选题
%% chapter: 曲线运动
%% point: 人造地球卫星
%% method: 无
%% score: 4
%% degree: 650
%% duplicateId: 0
%% body:
“天舟一号”货运飞船于 2017 年 4 月 20 日在文昌航天发射中心成功发射升空，与“天宫二号”空间实验室对接前，“天舟一号”在距离地面约 $380\Ukm$ 的圆轨道上飞行，则其\xzanswer{BCD}

\fourchoices[answer=BCD]
{角速度小于地球自转角速度}
{线速度小于第一宇宙速度}
{周期小于地球自转周期}
{向心加速度小于地面的重力加速度}

%% answer: BCD
%% memo:
\memoanswer{
“天舟一号”在距地面约 $380\Ukm$ 的圆轨道上飞行时，“天舟一号”的轨道半径远小于同步卫星的轨道半径，由 $G\frac{Mm}{r^{2}} = mr\omega^{2}$ 可知，其角速度大于同步卫星的角速度，即大于地球自转的角速度，故 A 错误；由于第一宇宙速度是卫星最大的环绕速度，因此“天舟一号”在轨的线速度小于第一宇宙速度，故 B 正确；由 $T = \frac{2\pi}{\omega}$ 可知，“天舟一号”做圆周运动的周期小于地球自转的周期，故 C 正确；由 $G\frac{Mm}{R^{2}} = mg$，$G\frac{Mm}{(R + h)^{2}} = ma$ 可知，向心加速度 $a$ 小于地球表面的重力加速度 $g$，故 D 正确。
}

\item
%% number: 7
%% paperName: 2017年江苏高考第 7 题 4 分
%% typeId: 2
%% type: 多选题
%% chapter: 交流电
%% point: LC振荡回路
%% method: 无
%% score: 4
%% degree: 850
%% duplicateId: 0
%% body:
某音响电路的简化电路图如图所示，输入信号既有高频成分，也有低频成分，则\xzanswer{BD}

\onepicture[w=5.6cm]{figs/07.png}

\fourchoices[answer=BD]
{电感 $L_{1}$ 的作用是通高频}
{电容 $C_{2}$ 的作用是通高频}
{扬声器甲用于输出高频成分}
{扬声器乙用于输出高频成分}

%% answer: BD
%% memo:
\memoanswer{
电感线圈有通低频、阻高频的作用，故 A 错误；电容器有通高频、阻低频的作用，故 B 正确；由此可以判断，$L_{1}$、$L_{2}$ 通低频，阻高频，$C_{1}$、$C_{2}$ 通高频，阻低频，因此扬声器甲用于输出低频成分，扬声器乙用于输出高频成分，故 C 错误 D 正确。
}

\item
%% number: 8
%% paperName: 2017年江苏高考第 8 题 4 分
%% typeId: 2
%% type: 多选题
%% chapter: 电场
%% point: 电场中的图象
%% method: 无
%% score: 4
%% degree: 550
%% duplicateId: 0
%% body:
在 $x$ 轴上有两个点电荷 $q_{1}$、$q_{2}$，其静电场的电势 $\varphi$ 在 $x$ 轴上分布如图所示。下列说法正确有\xzanswer{AC}

\onepicture[w=5.2cm]{\includesvg{figs/08}}

\fourchoices[answer=AC]
{$q_{1}$ 和 $q_{2}$ 带有异种电荷}
{$x_{1}$ 处的电场强度为零}
{负电荷从 $x_{1}$ 移到 $x_{2}$，电势能减小}
{负电荷从 $x_{1}$ 移到 $x_{2}$，受到的电场力增大}

%% answer: AC
%% memo:
\memoanswer{
由图象可知，空间的电势分布有正有负，则两个点电荷必定为异种电荷，故 A 正确；由 $E = \frac{\Delta\varphi}{\Delta x}$ 可知，图象的切线斜率为电场强度，因此 $x_{1}$ 处的电场强度不为零，故 B 错误；负电荷从 $x_{1}$ 移到 $x_{2}$ 的过程中，电势升高，电场强度减小，由 $E_{p} = q\varphi$，$F = qE$ 可知，电势能减小，受到的电场力减小，故 C 正确 D 错误。
}

\item
%% number: 9
%% paperName: 2017年江苏高考第 9 题 4 分
%% typeId: 2
%% type: 多选题
%% chapter: 功和能
%% point: 动能定理
%% method: 无
%% score: 4
%% degree: 350
%% duplicateId: 0
%% body:
如图所示，三个小球 A、B、C 的质量均为 $m$，A 与 B、C 间通过铰链用轻杆连接，杆长为 $L$，B、C 置于水平地面上，用一轻质弹簧连接，弹簧处于原长。现 A 由静止释放下降到最低点，两轻杆间夹角 $\alpha$ 由 $60^{\circ}$ 变为 $120^{\circ}$，A、B、C 在同一竖直平面内运动，弹簧在弹性限度内，忽略一切摩擦，重力加速度为 $g$。则此下降过程中\xzanswer{AB}

\onepicture[w=5.6cm]{\includesvg{figs/09}}

\fourchoices[answer=AB]
{A 的动能达到最大前，B 受到地面的支持力小于 $\frac{3}{2}mg$}
{A 的动能最大时，B 受到地面的支持力等于 $\frac{3}{2}mg$}
{弹簧的弹性势能最大时，A 的加速度方向竖直向下}
{弹簧的弹性势能最大值为 $\frac{\sqrt 3}{2}mgL$}

%% answer: AB
%% memo:
\memoanswer{
当 A 的动能达到最大前，A 向下加速运动，此时 A 处于失重状态，则整个系统对地面的压力小于 $3mg$，由对称性可知，B 对地面的压力小于 $\frac{3}{2}mg$，由牛顿第三定律可知，地面对 B 支持力小于 $\frac{3}{2}mg$，故 A 正确；当 A 的动能最大时，A 的加速度为零，这时系统既不失重，也不超重，系统对地面的压力等于 $3mg$，即 B 受到地面的支持力等于 $\frac{3}{2}mg$，故 B 正确；当弹簧的弹性势能最大时，A 减速运动到最低点，此时 A 的加速度方向竖直向上，故 C 错误；由机械能守恒定律可知，弹簧的弹性势能最大值等于 A 球的重力势能的减少量，即为 $mg(L\cos30^{\circ} - L\cos60^{\circ}) = \frac{\sqrt 3 - 1}{2}mgL$，故 D 错误。
}

\item
%% number: 10
%% paperName: 2017年江苏高考第 10 题 8 分
%% typeId: 4
%% type: 实验题
%% chapter: 功和能
%% point: 动能定理
%% method: 无
%% score: 8
%% degree: 550
%% duplicateId: 0
%% body:
利用图 \ref{2017江苏10a} 所示的实验装置探究恒力做功与物体动能变化的关系。小车的质量为 $M = 200.0\Ug$，钩码的质量为 $m = 10.0\Ug$，打点计时器的电源为 $50\UHz$ 的交流电。

\onepicture[num=1, numstyle=arabic, labelprefix={图}, label=2017江苏10a, w=6.4cm, align=right]{figs/10a.png}

\begin{enumerate}
	\item
	挂钩码前，为了消除摩擦力的影响，应调节木板右侧的高度，直至向左轻推小车观察到\tkanswer[2.5cm]{小车做匀速运动}。
	\item
	挂上钩码，按实验要求打出的一条纸带如图 \ref{2017江苏10b} 所示。选择某一点为 O，一次每隔 4 个计时点取一个计数点。用刻度尺量出相邻计数点间的距离 $\Delta x$，记录在纸带上。计算打出各计数点时小车的速度 $v$，其中打出计数点“1”时小车的速度 $v_{1} = \tkanswer[1.5cm]{0.228}\Ums$。

	\onepicture[num=2, numstyle=arabic, labelprefix={图}, label=2017江苏10b, w=9.5cm]{figs/10b.png}

	\item
	将钩码的重力视为小车受到的拉力，取 $g = 9.80\Umsq$，利用 $W = mg\Delta x$ 算出拉力对小车做的功 $W$。利用 $E_{k} = \frac{1}{2}Mv^{2}$ 算出小车动能，并求出动能的变化量 $\Delta E_{k}$。计算结果见下表。

	\begin{center}
	\begin{tabular}{|c|c|c|c|c|c|}
	\hline
	$W/\times10^{-3}\UJ$ & 2.45 & 2.92 & 3.35 & 3.81 & 4.26 \\
	\hline
	$\Delta E_{k}/\times10^{-3}\UJ$ & 2.31 & 2.73 & 3.12 & 3.61 & 4.00 \\
	\hline
	\end{tabular}
	\end{center}

	请根据表中的数据，在答题卡的方格纸上作出 $\Delta E_{k}-W$ 图象。

	\onepicture[w=7.0cm]{\includesvg{figs/10c}}

	\item
	实验结果表明，$\Delta E_{k}$ 总是略小于 $W$。某同学猜想是由于小车所受拉力小于钩码重力造成的。用题中小车和钩码质量的数据可算出小车受到的实际拉力 $F = \tkanswer[1.5cm]{0.093}\UN$。
\end{enumerate}

%% answer:
\jdanswer{
\begin{enumerate}
	\item
	小车做匀速运动
	\item
	0.228
	\item
	（作图见解析）
	\item
	0.093
\end{enumerate}
}
%% memo:
\memoanswer{
（1）挂钩码前，调节木板右侧的高度，直至向左轻推小车时，观察到小车做匀速直线运动（或观察纸带，打出的点间隔均匀），此时消除了摩擦力的影响。

（2）打计数点“1”时小车的速度 $v_{1} = \frac{(2.06 + 2.50)\times10^{-2}}{0.2}\Ums = 0.228\Ums$。

（3）由表格中的数据，描点作图如答图所示。

\onepicture[w=7.0cm]{figs/10_an.png}

（4）由功的定义得 $W = mgx$，由动能定理得 $\Delta E_{k} = Fx$，则 $\Delta E_{k} = \frac{F}{mg}W$，结合图象有 $F = \frac{4.0 - 2.0}{4.26 - 2.15}\times0.098\UN = 0.093\UN$。
}

\item
%% number: 11
%% paperName: 2017年江苏高考第 11 题 10 分
%% typeId: 4
%% type: 实验题
%% chapter: 电路
%% point: 传感器
%% method: 无
%% score: 10
%% degree: 450
%% duplicateId: 0
%% body:
某同学通过实验制作一个简易的温控装置，实验原理电路图如图 \ref{2017江苏11a} 所示，继电器与热敏电阻 $R_{1}$、滑动变阻器 $R$ 串联接在电源 $E$ 两端，当继电器的电流超过 $15\UmA$ 时，衔铁被吸合，加热器停止加热，实现温控。继电器的电阻约为 $20\UO$，热敏电阻的阻值 $R_{t}$ 与温度 $t$ 的关系如下表所示：

\begin{center}
\begin{tabular}{|c|c|c|c|c|c|c|}
\hline
$t/\UdC$ & 30.0 & 40.0 & 50.0 & 60.0 & 70.0 & 80.0 \\
\hline
$R_{t}/\UO$ & 199.5 & 145.4 & 108.1 & 81.8 & 62.9 & 49.1 \\
\hline
\end{tabular}
\end{center}

\onepicture[num=1, numstyle=arabic, labelprefix={图}, label=2017江苏11a, w=8.8cm]{figs/11a.png}

\begin{enumerate}
	\item
	提供的实验器材有：电源 $E_{1}$（$3\UV$，内阻不计）、电源 $E_{2}$（$6\UV$，内阻不计）、滑动变阻器 $R_{1}$（$0\sim200\UO$）、滑动变阻器 $R_{2}$（$0\sim500\UO$）、热敏电阻 $R_{t}$，继电器、电阻箱（$0\sim999.9\UO$）、开关 S、导线若干。

	为使该装置实现对 $30\sim80\UdC$ 之间任一温度的控制，电源 $E$ 应选用\tkanswer[1.2cm]{$E_{2}$}（选填“$E_{1}$”或“$E_{2}$”），滑动变阻器 $R$ 应选用\tkanswer[1.2cm]{$R_{2}$}（选填“$R_{1}$”或“$R_{2}$”）。
	\item
	实验发现电路不工作。某同学为排查电路故障，用多用电表测量各接点间的电压，则应将如图 \ref{2017江苏11b} 所示的选择开关旋至\tkanswer[0.8cm]{C}（选填“A”、“B”、“C”或“D”）。
	\item
	合上开关 S，用调节好的多用电表进行排查，在图 \ref{2017江苏11a} 中，若只有 b、c 间断路，则应发现表笔接入 a、b 时指针\tkanswer[2cm]{不偏转}（选填“偏转”或“不偏转”），接入 a、c 时指针\tkanswer[2cm]{偏转}（选填“偏转”或“不偏转”）。
	\item
	排除故障后，欲使衔铁在热敏电阻为 $50\UdC$ 时被吸合，下列操作步骤正确顺序是\tkanswer[3cm]{⑤④②③①}。（填写各步骤前的序号）

	\begin{steps}
		\item
		将热敏电阻接入电路
		\item
		观察到继电器的衔铁被吸合
		\item
		断开开关，将电阻箱从电路中移除
		\item
		合上开关，调节滑动变阻器的阻值
		\item
		断开开关，用电阻箱替换热敏电阻，将阻值调至 $108.1\UO$
	\end{steps}
\end{enumerate}

\onepicture[num=2, numstyle=arabic, labelprefix={图}, label=2017江苏11b, w=5.2cm]{figs/11b.png}

%% answer:
\jdanswer{
\begin{enumerate}
	\item
	$E_{2}$，$R_{2}$
	\item
	C
	\item
	不偏转，偏转
	\item
	⑤④②③①
\end{enumerate}
}
%% memo:
\memoanswer{
（1）如果电源用 $E_{1}$，则在 $t = 30\UdC$ 时电路中的最大电流 $I_{m} = \frac{3}{199.5 + 20}\UA = 13.67\UmA < 15\UmA$，不能实现对此温度的控制，因此电源应用 $E_{2}$；为了实现在 $t = 80\UdC$ 时对电路的控制，设滑动变阻器的最大值至少为 $R$，则 $\frac{6\UV}{R + 20\UO + 49.1\UO} = 0.015\UA$，求得 $R = 330.9\UO$，因此滑动变阻器应选用 $R_{2}$。

（2）要用多用电表的直流电压挡探测故障，应将选择开关旋至面板直流电压挡 C 上。

（3）如果只有 b、c 间断路，说明 b 点与电源的负极间不通，a、b 间的电压为零，表笔接在 a、b 间时，指针不偏转；c 点与电源的负极相通，a 与电源的正极相通，a、c 间有电压，因此两表笔接入 a、c 间时指针发生偏转。

（4）排除故障后，欲使衔铁在热敏电阻为 $50\UdC$ 时被吸合，应先断开开关，用电阻箱替换热敏电阻，并将阻值调至 $108.1\UO$，合上开关，调节滑动变阻器的阻值，直至观察到继电器的衔铁被吸合，这时断开开关，将电阻箱从电路中移除，将热敏电阻再接入电路。因此操作步骤的正确顺序是⑤④②③①。
}

\item
%% number: 12
%% paperName: 2017年江苏高考第 12 题 4 分
%% typeId: 6
%% type: 计算题
%% chapter: 运动和力
%% point: 动量守恒定律
%% method: 无
%% score: 4
%% degree: 750
%% duplicateId: 0
%% body:
【选做题】本题包括 A、B、C 三小题，请选定其中两小题，并在相应的答题区域内作答。若多做，则按 A、B 两小题评分。

\textbf{A．}（选修 3-3）（12 分）

\begin{enumerate}
	\item
	一定质量的理想气体从状态 A 经过状态 B 变化到状态 C，其 $V-T$ 图象如图所示。下列说法正确的有\xzanswer{BC}

	\onepicture[w=4.0cm, align=right]{\includesvg{figs/12a}}

	\fourchoices[answer=BC]
	{$A\to B$ 的过程中，气体对外界做功}
	{$A\to B$ 的过程中，气体放出热量}
	{$B\to C$ 的过程中，气体压强不变}
	{$A\to B\to C$ 的过程中，气体内能增加}
	\item
	（甲）和（乙）图中是某同学从资料中查到的两张记录水中炭粒运动位置连线的图片，记录炭粒位置的时间间隔均为 $30\Us$，两方格纸每格表示的长度相同。比较两张图片可知：若水温相同，\tkanswer[1cm]{甲}（选填“甲”或“乙”）中炭粒的颗粒较大；若炭粒大小相同，\tkanswer[1cm]{乙}（选填“甲”或“乙”）中水分子的热运动较剧烈。

	\twopicture[numstyle=chinese, labelprefix={}, labelA=2017江苏12b, labelB=2017江苏12c, wA=3.4cm, wB=3.4cm]{figs/12b.png}{figs/12c.png}
	\item
	科学家可以运用无规则运动的规律来研究生物蛋白分子。资料显示，某种蛋白的摩尔质量为 $66\Ukgmol$，其分子可视为半径为 $3\times10^{-9}\Um$ 的球，已知阿伏伽德罗常数为 $6.0\times10^{23}\Umoln$。请估算该蛋白的密度。（计算结果保留一位有效数字）
\end{enumerate}

\textbf{B．}（选修 3-4）（12 分）

\begin{enumerate}
	\item
	接近光速飞行的飞船和地球上各有一只相同的铯原子钟，飞船和地球上的人观测这两只钟的快慢，下列说法正确的有\xzanswer{AC}

	\fourchoices[answer=AC]
	{飞船上的人观测到飞船上的钟较快}
	{飞船上的人观测到飞船上的钟较慢}
	{地球上的人观测到地球上的钟较快}
	{地球上的人观测到地球上的钟较慢}
	\item
	野生大象群也有自己的“语言”。研究人员录下象群“语言”交流时发出的声音，发现以 2 倍速度快速播放时，能听到比正常播放时更多的声音。播放速度变为原来的 2 倍时，播出声波的\tkanswer[1.5cm]{频率}（选填“周期”或“频率”）也变为原来的 2 倍，声波的传播速度\tkanswer[1.5cm]{不变}（选填“变大”、“变小”或“不变”）。
	\item
	人的眼球可简化为如图所示的模型，折射率相同、半径不同的两个球体共轴，平行光束宽度为 $D$，对称地沿轴线方向射入半径为 $R$ 的小球，会聚在轴线上的 P 点。取球体的折射率为 $\sqrt 2$，且 $D = \sqrt 2 R$，求光线的会聚角 $\alpha$。（示意图未按比例画出）

	\onepicture[w=5.5cm, align=right]{figs/12d.png}
\end{enumerate}

\textbf{C．}（选修 3-5）（12 分）

\begin{enumerate}
	\item
	原子核的比结合能曲线如图所示，根据该曲线，下列判断中正确的有\xzanswer{BC}

	\onepicture[w=5.5cm, align=right]{figs/12e.png}

	\fourchoices[answer=BC]
	{\ce{^{4}_{2}He} 核的结合能约为 $14\UMeV$}
	{\ce{^{4}_{2}He} 核比 \ce{^{6}_{3}Li} 核更稳定}
	{两个 \ce{^{2}_{1}H} 核结合成 \ce{^{4}_{2}He} 核时释放能量}
	{\ce{^{235}_{92}U} 核中核子的平均结合能比 \ce{^{89}_{36}Kr} 核中的大}
	\item
	质子和 $\alpha$ 粒子被加速到相同动能时，质子的动量\tkanswer[1.2cm]{小于}（选填“大于”、“小于”或“等于”）$\alpha$ 粒子的动量，质子和 $\alpha$ 粒子的德布罗意波波长之比为\tkanswer[1.5cm]{$2:1$}。
	\item
	甲、乙两运动员在做花样滑冰表演，沿同一直线相向运动，速度大小都是 $1\Ums$，甲、乙相遇时用力推对方，此后都沿各自原方向的反方向运动，速度大小分别为 $1\Ums$ 和 $2\Ums$。求甲、乙两运动员的质量之比。
\end{enumerate}

%% answer:
\jdanswer{
\begin{enumerate}
	\item
	A：（1）BC；（2）甲，乙；（3）$\rho = 1\times10^{3}\Ukgmc$（或 $\rho = 5\times10^{2}\Ukgmc$，$5\times10^{2}\sim5\times10^{3}$ 都算对）。
	\item
	B：（1）AC；（2）频率，不变；（3）$\alpha = 30^{\circ}$。
	\item
	C：（1）BC；（2）小于，$2:1$；（3）$\frac{3}{2}$。
\end{enumerate}
}
%% memo:
\memoanswer{
\textbf{A．}（1）由图可知，从状态 A 到状态 B 气体的体积减小，外界对气体做功，故 A 错误；从状态 A 到状态 B 过程气体的温度不变，内能不变，由热力学第一定律可知，气体放出热量，故 B 正确；由理想气体状态方程 $pV = nRT$，$V = \frac{nR}{p}T$ 可知，从状态 B 到状态 C 气体发生的是等压变化，气体的温度在降低，内能在减少，故 C 正确 D 错误。

（2）影响布朗运动快慢的因素有两个，悬浮颗粒的大小和温度，炭粒越小、温度越高布朗运动越明显，从图可以看出，乙图反映的布朗运动明显，因此温度相同时，甲图中的炭粒大；炭粒大小相同时，乙图中水的温度高。

（3）摩尔体积 $V = \frac{4}{3}\pi r^{3}N_{A}$（或 $V = (2r)^{3}N_{A}$），由密度公式知 $\rho = \frac{M}{V}$，联立解得 $\rho = \frac{3M}{4\pi r^{3}N_{A}}$（或 $\rho = \frac{M}{8r^{3}N_{A}}$），代入数据得 $\rho = 1\times10^{3}\Ukgmc$（或 $\rho = 5\times10^{2}\Ukgmc$，$5\times10^{2}\sim5\times10^{3}\Ukgmc$ 都算对）。

\textbf{B．}（1）相对论告诉我们，运动的钟会变慢，由于飞船上的人观测飞船上的钟是相对静止的，观测到地球上的钟是相对运动的，因此飞船上的人观察到飞船上的钟比地球上的钟快，故 A 正确 B 错误；同样，地球上的人观测到飞船上的钟是相对运动的，观测到地球上的钟是相对静止的，因此他观测到地球上的钟比飞船上的钟快，故 C 正确 D 错误。

（2）播放速度变为原来的 2 倍，则播出声波的周期为原来的一半，频率为原来的 2 倍，但在同一介质中声波传播速度不变。

（3）由几何关系得 $\sin i = \frac{D}{2R}$，解得 $i = 45^{\circ}$，由折射定律得 $n = \frac{\sin i}{\sin\gamma}$，解得 $\gamma = 30^{\circ}$，且 $i = \gamma + \frac{\alpha}{2}$，解得 $\alpha = 30^{\circ}$。

\textbf{C．}（1）由图可知，\ce{^{4}_{2}He} 的比结合能约为 $7\UMeV$，因此 \ce{^{4}_{2}He} 的结合能约为 $7\UMeV\times4 = 28\UMeV$，故 A 错误；比结合能越大，原子核越牢固，结合图象可知，故 B 正确；两个比结合能小的 \ce{^{2}_{1}H} 核结合成比结合能大的 \ce{^{4}_{2}He}，会释放能量，故 C 正确；由图象可知，\ce{^{235}_{92}U} 的比结合能（即平均结合能）比 \ce{^{89}_{36}Kr} 的小，故 D 错误。

（2）由 $p = \sqrt{2mE_{k}}$ 可知，动能相同时，质子的质量比 $\alpha$ 粒子的质量小，因此动量小；由 $\lambda = \frac{h}{p}$ 可知，质子和 $\alpha$ 粒子的德布罗意波波长之比为 $\frac{\lambda_{\text{p}}}{\lambda_{\alpha}} = \sqrt{\frac{m_{\alpha}}{m_{\text{p}}}} = \sqrt{\frac{4}{1}} = \frac{2}{1}$。

（3）由动量守恒得 $m_{1}v_{1} - m_{2}v_{2} = m_{2}v_{2}' - m_{1}v_{1}'$，解得 $\frac{m_{1}}{m_{2}} = \frac{v_{2} + v_{2}'}{v_{1} + v_{1}'}$，代入数据得 $\frac{m_{1}}{m_{2}} = \frac{3}{2}$。
}

\item
%% number: 13
%% paperName: 2017年江苏高考第 13 题 15 分
%% typeId: 6
%% type: 计算题
%% chapter: 电磁感应
%% point: 电磁感应受力分析
%% method: 无
%% score: 15
%% degree: 550
%% duplicateId: 0
%% body:
如图所示，两条相距 $d$ 的平行金属导轨位于同一水平面内，其右端接一阻值为 $R$ 的电阻。质量为 $m$ 的金属杆静置在导轨上，其左侧的矩形匀强磁场区域 MNPQ 的磁感应强度大小为 $B$、方向竖直向下。当该磁场区域以速度 $v_{0}$ 匀速地向右扫过金属杆后，金属杆的速度变为 $v$。导轨和金属杆的电阻不计，导轨光滑且足够长，杆在运动过程中始终与导轨垂直且两端与导轨保持良好接触。求：

\begin{enumerate}
	\item
	MN 刚扫过金属杆时，杆中感应电流的大小 $I$；
	\item
	MN 刚扫过金属杆时，杆的加速度大小 $a$；
	\item
	PQ 刚要离开金属杆时，感应电流的功率 $P$。
\end{enumerate}

\onepicture[w=7.0cm, align=right]{\includesvg{figs/13}}

%% answer:
\jdanswer{
\begin{enumerate}
	\item
	$I = \frac{Bdv_{0}}{R}$
	\item
	$a = \frac{B^{2}d^{2}v_{0}}{mR}$
	\item
	$P = \frac{B^{2}d^{2}(v_{0} - v)^{2}}{R}$
\end{enumerate}
}
%% memo:
\memoanswer{
（1）由法拉第电磁感应定律知，感应电动势为 $E = Bdv_{0}$，感应电流为 $I = \frac{E}{R}$，联立解得 $I = \frac{Bdv_{0}}{R}$。

（2）金属杆受到的安培力大小为 $F = BId$，由牛顿第二定律得 $F = ma$，联立解得 $a = \frac{B^{2}d^{2}v_{0}}{mR}$。

（3）金属杆切割磁感线的速度大小为 $v' = v_{0} - v$，感应电动势大小为 $E = Bd(v_{0} - v)$，电功率大小为 $P = \frac{E^{2}}{R}$，联立解得 $P = \frac{B^{2}d^{2}(v_{0} - v)^{2}}{R}$。
}

\item
%% number: 14
%% paperName: 2017年江苏高考第 14 题 16 分
%% typeId: 6
%% type: 计算题
%% chapter: 力物体的平衡
%% point: 物体系平衡
%% method: 无
%% score: 16
%% degree: 450
%% duplicateId: 0
%% body:
如图所示，两个半圆柱 A、B 紧靠着静置于水平地面上，其上有一光滑圆柱 C，三者半径均为 $R$。C 的质量为 $m$，A、B 的质量都为 $\frac{m}{2}$，与地面的动摩擦因数均为 $\mu$。现用水平向右的力拉 A，使 A 缓慢移动，直至 C 恰好降到地面。整个过程中 B 保持静止。设最大静摩擦力等于滑动摩擦力，重力加速度为 $g$。求：

\begin{enumerate}
	\item
	未拉 A 时，C 受到 B 作用力的大小 $F$；
	\item
	动摩擦因数的最小值 $\mu_{\min}$；
	\item
	A 移动的整个过程中，拉力做的功 $W$。
\end{enumerate}

\onepicture[w=6.4cm, align=right]{\includesvg{figs/14}}

%% answer:
\jdanswer{
\begin{enumerate}
	\item
	$F = \frac{\sqrt 3}{3}mg$
	\item
	$\mu_{\min} = \frac{\sqrt 3}{2}$
	\item
	$W = (2\mu - 1)(\sqrt 3 - 1)mgR$
\end{enumerate}
}
%% memo:
\memoanswer{
（1）由 C 受力平衡得 $2F\cos30^{\circ} = mg$，解得 $F = \frac{\sqrt 3}{3}mg$。

（2）当 C 恰好降到地面时，B 受 C 压力的水平分力最大，大小为 $F_{x\max} = \frac{\sqrt 3}{2}mg$，B 受地面的摩擦力大小为 $f = \mu mg$，由题意知 $f_{\min} = F_{x\max}$，联立解得 $\mu_{\min} = \frac{\sqrt 3}{2}$。

（3）A 移动时，C 下降的高度为 $h = (\sqrt 3 - 1)R$，A 的位移大小为 $x = 2(\sqrt 3 - 1)R$，摩擦力做功的大小为 $W_{f} = fx = 2(\sqrt 3 - 1)\mu mgR$，由动能定理得 $W - W_{f} + mgh = 0 - 0$，解得 $W = (2\mu - 1)(\sqrt 3 - 1)mgR$。
}

\item
%% number: 15
%% paperName: 2017年江苏高考第 15 题 16 分
%% typeId: 6
%% type: 计算题
%% chapter: 磁场
%% point: 带电粒子在磁场中的运动
%% method: 无
%% score: 16
%% degree: 250
%% duplicateId: 0
%% body:
一台质谱仪的工作原理如图所示。大量的甲、乙两种离子飘入电压为 $U_{0}$ 的加速电场，其初速度几乎为 0，经过加速后，通过宽为 $L$ 的狭缝 MN 沿着与磁场垂直的方向进入磁感应强度为 $B$ 的匀强磁场中，最后打到照相底片上。已知甲、乙两种离子的电荷量均为 $+q$，质量分别为 $2m$ 和 $m$，图中虚线为经过狭缝左、右边界 M、N 的甲种离子的运动轨迹。不考虑离子间的相互作用。

\begin{enumerate}
	\item
	求甲种离子打在底片上的位置到 N 点的最小距离 $x$；
	\item
	在图中用斜线标出磁场中甲种离子经过的区域，并求该区域最窄处的宽度 $d$；
	\item
	若考虑加速电压有波动，在 $(U_{0} - \Delta U)$ 到 $(U_{0} + \Delta U)$ 之间变化，要使甲、乙两种离子在底片上没有重叠，求狭缝宽度 $L$ 满足的条件。
\end{enumerate}

\onepicture[w=6.0cm, align=right]{\includesvg{figs/15}}

%% answer:
\jdanswer{
\begin{enumerate}
	\item
	$x = \frac{4}{B}\sqrt{\frac{mU_{0}}{q}} - L$
	\item
	$d = \frac{2}{B}\sqrt{\frac{mU_{0}}{q}} - \sqrt{\frac{4mU_{0}}{qB^{2}} - \frac{L^{2}}{4}}$
	\item
	$L < \frac{2}{B}\sqrt{\frac{m}{q}}\left[2\sqrt{U_{0} - \Delta U} - \sqrt{2(U_{0} + \Delta U)}\right]$
\end{enumerate}
}
%% memo:
\memoanswer{
（1）设甲种离子在磁场中的运动半径为 $r_{1}$。电场加速时，由功能关系得 $qU_{0} = \frac{1}{2}\times2mv^{2}$，由洛伦兹力提供向心力得 $qvB = 2m\frac{v^{2}}{r_{1}}$，联立解得 $r_{1} = \frac{2}{B}\sqrt{\frac{mU_{0}}{q}}$，由几何关系得 $x = 2r_{1} - L$，解得 $x = \frac{4}{B}\sqrt{\frac{mU_{0}}{q}} - L$。

（2）离子经过的区域见下图。最窄处位于两虚线交点的垂线上，由几何关系得 $d = r_{1} - \sqrt{r_{1}^{2} - \left(\frac{L}{2}\right)^{2}}$，解得 $d = \frac{2}{B}\sqrt{\frac{mU_{0}}{q}} - \sqrt{\frac{4mU_{0}}{qB^{2}} - \frac{L^{2}}{4}}$。

\onepicture[w=7.5cm]{figs/15_an.png}

（3）设乙种离子在磁场中的运动半径为 $r_{2}$。$r_{1}$ 的最小半径为 $r_{1\min} = \frac{2}{B}\sqrt{\frac{m(U_{0} - \Delta U)}{q}}$，$r_{2}$ 的最大半径为 $r_{2\max} = \frac{1}{B}\sqrt{\frac{2m(U_{0} + \Delta U)}{q}}$，由题意知 $2r_{1\min} - 2r_{2\max} > L$，即 $\frac{4}{B}\sqrt{\frac{m(U_{0} - \Delta U)}{q}} - \frac{2}{B}\sqrt{\frac{2m(U_{0} + \Delta U)}{q}} > L$，解得 $L < \frac{2}{B}\sqrt{\frac{m}{q}}\left[2\sqrt{U_{0} - \Delta U} - \sqrt{2(U_{0} + \Delta U)}\right]$。
}

\end{enumerate}

\end{document}
